\section*{अध्याय \ref{ch:b2:mantle-melting} --- प्रावार का पिघलना}

\begin{solution}{exo:b2:mantle-melting:1}
(क) \qty{90}{\percent} से अधिक ओलिवीन: \omterm{def:b2:mantle-melting:peridotites}{ड्यूनाइट}। (ख) ओलिवीन और ऑर्थोपाइरॉक्सीन,
बिना क्लाइनोपाइरॉक्सीन के: \omterm{def:b2:mantle-melting:peridotites}{हार्ज़बर्गाइट}। (ग) दोनों पाइरॉक्सीनों के साथ \qty{40}{\percent} से अधिक
ओलिवीन: \omterm{def:b2:mantle-melting:peridotites}{लर्ज़ोलाइट}। (घ) \qty{40}{\percent} से कम ओलिवीन: एक पाइरॉक्सिनाइट
(एक ओलिवीन-वेबस्टराइट, दोनों पाइरॉक्सीनों के साथ)। (ङ) $57/98 = \qty{58.2}{\percent}$,
$28/98 = \qty{28.6}{\percent}$, $13/98 = \qty{13.3}{\percent}$: एक \omterm{def:b2:mantle-melting:peridotites}{लर्ज़ोलाइट}।
\end{solution}

\begin{solution}{exo:b2:mantle-melting:2}
(क) \qty{60}{\percent} एनॉर्थाइट पर लिक्विडस \qty{1390}{\celsius} है,
एनॉर्थाइट क्षेत्र में: एनॉर्थाइट पहले क्रिस्टलित होता है। (ख) डायोप्साइड
गलनक्रांतिक पर, \qty{1270}{\celsius} पर, प्रकट होता है, जब द्रव में \qty{42}{\percent} एनॉर्थाइट होता है।
(ग) दोनों खनिजों के किसी भी मिश्रण का पहला द्रव गलनक्रांतिक द्रव है,
\qty{42}{\percent} एनॉर्थाइट, \qty{1270}{\celsius} पर।
\end{solution}

\begin{solution}{exo:b2:mantle-melting:3}
(क) $\num{4.0e-5}\times1623/(3300\times1000) = \qty{1.97e-8}{K/Pa} =
\qty{19.7}{K/GPa}$। (ख) $\num{e9}/(3300\times9.81) = \qty{3.09e4}{m}$, या
\qty{30.9}{km} प्रति गीगापास्कल; $19.7/30.9 = \qty{0.64}{K/km}$, पृष्ठ के निकट के
\qty{30}{K/km} से लगभग पचास गुना कम, जहाँ ऊष्मा चालित होती है, वहन नहीं
की जाती।
\end{solution}

\begin{solution}{exo:b2:mantle-melting:4}
(क) $T_{\mathrm s}(0) = \qty{1085.7}{\celsius}$; $T_{\mathrm s}(1) = 1085.7 +
132.9 - 5.1 = \qty{1213.5}{\celsius}$; $T_{\mathrm s}(2) = 1085.7 + 265.8 - 20.4 =
\qty{1331.1}{\celsius}$; $T_{\mathrm s}(3) = 1085.7 + 398.7 - 45.9 =
\qty{1438.5}{\celsius}$। (ख) $132.9 - 2\times5.1\times2 = \qty{112.5}{K/GPa}$।
(ग) \qty{3}{GPa} पर \qty{1400}{\celsius} सॉलिडस,
\qty{1438.5}{\celsius}, से नीचे है; रुद्धोष्म के अनुदिश \qty{2}{GPa} तक उठते हुए वह केवल
लगभग \qty{20}{\celsius} ठंडा होता है, लगभग \qty{1380}{\celsius} तक, जबकि सॉलिडस गिरकर
\qty{1331.1}{\celsius} हो जाता है: चट्टान अब अपने सॉलिडस से ऊपर है और उसका एक भाग पिघलता है।
\end{solution}

\begin{solution}{exo:b2:mantle-melting:5}
(क) $1350 + (1277.5 - 1350)\times(30 - 20)/(40 - 20) = \qty{1314}{\celsius}$।
(ख) संघटन गलनक्रांतिक के डायोप्साइड वाले पक्ष पर है: डायोप्साइड।
(ग) द्रव \qty{42}{\percent} एनॉर्थाइट के साथ गलनक्रांतिक तक पहुँचता है; एनॉर्थाइट
द्रव में रहता है, इसलिए द्रव-अंश $30/42 = 0.714$ है और तंत्र का $1 - 0.714 =
\qty{29}{\percent}$ डायोप्साइड के रूप में क्रिस्टलित हो चुका है। (घ) सब कुछ
क्रिस्टलित होता है: \qty{70}{\percent} डायोप्साइड और \qty{30}{\percent} एनॉर्थाइट,
समग्र संघटन।
\end{solution}

\begin{solution}{exo:b2:mantle-melting:6}
(क) फ़ॉर्स्टेराइट: $60.08/(2\times40.30 + 60.08) = 60.08/140.68 =
\qty{42.71}{\percent}$ \ce{SiO2}; एन्स्टेटाइट: $60.08/100.38 = \qty{59.85}{\percent}$।
(ख) $0.975\times59.85 + 2.5 = \qty{60.85}{\percent}$। (ग) एन्स्टेटाइट फ़ॉर्स्टेराइट
और द्रव के बीच है: फ़ॉर्स्टेराइट का अंश $(60.85 - 59.85)/(60.85 -
42.71) = 1.00/18.14 = \qty{5.5}{\percent}$ है, बोवेन और एंडरसन का मान। (घ)
गलित उस ऑर्थोपाइरॉक्सीन से अधिक सिलिका-समृद्ध है जिससे वह आता है, और अवशेष
ओलिवीन पाता है।
\end{solution}

\begin{solution}{exo:b2:mantle-melting:7}
(क) $5.1P^2 - 112.9P + 214.3 = 0$: $P_0 = (112.9 - \sqrt{12746.4 -
4371.7})/10.2 = (112.9 - 91.51)/10.2 = \qty{2.10}{GPa}$, जो $2.097\times30.9 =
\qty{64.8}{km}$ पर है। (ख) $5.1P^2 - 112.9P + 364.3 = 0$: $P_0 = (112.9 -
72.90)/10.2 = \qty{3.92}{GPa}$, \qty{121}{km} पर। (ग) $T_{\mathrm s}'(2.10) = 132.9
- 21.4 = \qty{111.5}{K/GPa}$, इसलिए $\Delta P_0 = 10/(111.5 - 20) = \qty{0.109}{GPa}$
प्रति \qty{10}{\celsius}, या \qty{3.4}{km}।
\end{solution}

\begin{solution}{exo:b2:mantle-melting:8}
(क) $T_{\mathrm s} = \qty{1331.1}{\celsius}$, $T_{\mathrm l} = 1475.0 + 160.0 -
12.8 = \qty{1622.2}{\celsius}$; $T = 1331.1 + 0.10\times291.1 =
\qty{1360.2}{\celsius}$, या \qty{1633}{K}। (ख) $(\partial T/\partial P)_F = 112.5
+ 0.10\times(67.2 - 112.5) = \qty{108.0}{K/GPa}$; $(\partial T/\partial F)_P =
T_{\mathrm l} - T_{\mathrm s} = \qty{291.1}{K}$। (ग) $c_p/T = 0.612$; प्रति गीगापास्कल $\alpha/\rho
= \qty{12.1}{J/(kg.K)}$ के साथ अंश $0.612\times108.0 - 12.1
= 54.0$ है और हर $300 + 0.612\times291.1 = 478.2$: $-\mathrm dF/\mathrm
dP = 0.113$, \qty{11.3}{\percent} प्रति गीगापास्कल। (घ) $54.0/178.2 = 0.303$,
\qty{30}{\percent} प्रति गीगापास्कल: गुप्त ऊष्मा उत्पादकता को 2.7 के
गुणक से घटा देती है।
\end{solution}

\begin{solution}{exo:b2:mantle-melting:9}
(क) $43\times1 = \qty{43}{\celsius}$; $43\times5^{0.75} = 43\times3.34 =
\qty{144}{\celsius}$; $43\times13^{0.75} = 43\times6.85 = \qty{294}{\celsius}$।
(ख) $1213.5 - 294 = \qty{919}{\celsius}$। (ग) \qty{125}{ppm}
\qty{0.0125}{\percent} है, इसलिए $X = 0.0125/0.01 = \qty{1.25}{\percent}$ और $\Delta T =
43\times1.25^{0.75} = 43\times1.18 = \qty{51}{\celsius}$।
\end{solution}

\begin{solution}{exo:b2:mantle-melting:10}
(क) आधार $2H$ चौड़ा है और प्रावार उसे $U$ पर पार करता है: कटक की प्रति एकांक लंबाई
और समय $2UH$। हर भुजा ऊँचाई $H$ तक फैली है और $U$ पर क्षैतिज रूप से पार की जाती है:
हर एक $UH$, कुल $2UH$। भुजाएँ $\theta$ पर होतीं तो आधार
$2H/\tan\theta$ चौड़ा होता और अंतर्वाह $2UH/\tan\theta$, जो बहिर्वाह के बराबर केवल तब होता जब
$\tan\theta = 1$: इस मॉडल में कोण द्रव्यमान-संरक्षण से तय होता है। (ख)
स्थिर उत्पादकता के साथ शीर्ष पर \omterm{def:b2:mantle-melting:fraction}{गलन-अंश} $F_{\max} = \Pi P_0 = \Pi\rho
g z_0$ है, और $H = z_0$: $h_c = F_{\max}z_0/2 = \Pi\rho g z_0^2/2$। (ग) हर वर्ष
\qty{50}{mm} नई प्लेट, कटक के प्रति किलोमीटर $\qty{0.05}{m}\times\qty{1000}{m}\times\qty{7100}{m} =
\qty{3.6e5}{m^3}$ भूपर्पटी; \qty{12}{\percent} माध्य \omterm{def:b2:mantle-melting:fraction}{गलन-अंश}
पर यह $\num{3.6e5}/0.12 = \qty{3.0e6}{m^3}$ प्रावार से आती है।
\end{solution}

\begin{solution}{exo:b2:mantle-melting:11}
(क) द्रव के एकांक द्रव्यमान से A का द्रव्यमान $t$ निकल जाने के बाद, द्रव में A का $0.6 -
t$, B का 0.2 और C का 0.2 है, द्रव्यमान $1 - t$ में: $\big((0.6 - t), 0.2,
0.2\big)/(1 - t)$, और B/C 1 पर रहता है। (ख) कोटेक्टिक पर, $(0.40 - 0.15s, 0.60 -
0.30s, 0.45s)$; B = C से $0.60 - 0.30s = 0.45s$, $s = 0.8$: $L_1$ में A 28, B
36, C 36 है। तब $0.2/(1 - t) = 0.36$, $t = 0.44$: तंत्र का \qty{44}{\percent}
A के रूप में क्रिस्टलित हो चुका है। (ग) सारा C द्रव में है: $f = 20/45 = 0.44$,
\qty{44}{\percent} अब भी द्रव। (घ) A: $0.60 - 0.444\times0.25 = 0.49$; B: $0.20 -
0.444\times0.30 = 0.07$; द्रव सहित, $49 + 7 + 44 = \qty{100}{\percent}$।
\end{solution}

\begin{solution}{exo:b2:mantle-melting:12}
(क) $X_{\mathrm{sat}}(1) = 12 + 1 = \qty{13}{\percent}$, $\Delta T =
\qty{294}{\celsius}$, संतृप्त सॉलिडस $1213.5 - 294 = \qty{919}{\celsius}$;
$X_{\mathrm{sat}}(2) = 12\times1.516 + 2 = \qty{20.2}{\percent}$, $\Delta T =
43\times20.2^{0.75} = 43\times9.53 = \qty{410}{\celsius}$, संतृप्त सॉलिडस $1331.1
- 410 = \qty{921}{\celsius}$। (ख) ऊँचे दाब पर गलित अधिक जल घोलता है,
और बड़ी गिरावट शुष्क सॉलिडस की \qty{118}{\celsius} वृद्धि को निरस्त कर देती है।
(ग) सॉलिडस को $1331.1 - 1100 = \qty{231}{\celsius}$ गिरना होगा: $X =
(231/43)^{1/0.75} = 5.37^{1.333} = \qty{9.4}{\percent}$ जल, संतृप्ति के
\qty{20.2}{\percent} से कम, इसलिए कील पिघल सकती है।
\end{solution}

\begin{solution}{pb:b2:mantle-melting:1}
\textbf{1.} $1350 + 273 = \qty{1623}{K}$।
\textbf{2.} $\num{4.0e-5}\times1623/(3300\times1000) = \qty{1.97e-8}{K/Pa} =
\qty{19.7}{K/GPa}$, जो \qty{20}{K/GPa} तक पूर्णांकित होता है।
\textbf{3.} $\num{e9}/(3300\times9.81) = \qty{30.9}{km}$; $20/30.9 =
\qty{0.65}{K/km}$।
\textbf{4.} $1350 + 20P = 1085.7 + 132.9P - 5.1P^2$, अर्थात $5.1P^2 - 112.9P +
264.3 = 0$: $P_0 = (112.9 - \sqrt{12746.4 - 5391.7})/10.2 = (112.9 - 85.76)/10.2 =
\qty{2.66}{GPa}$।
\textbf{5.} $2.66\times30.9 = \qty{82}{km}$; $T = 1350 + 20\times2.66 =
\qty{1403}{\celsius}$।
\textbf{6.} $1085.7 + 132.9\times2.66 - 5.1\times2.66^2 = 1085.7 + 353.5 - 36.1 =
\qty{1403}{\celsius}$।
\textbf{7.} $T_{\mathrm s} = 1085.7 + 199.4 - 11.5 = \qty{1273.6}{\celsius}$;
$T_{\mathrm l} = 1475.0 + 120.0 - 7.2 = \qty{1587.8}{\celsius}$; अंतर
\qty{314.2}{K}।
\textbf{8.} $1273.6 + 0.10\times314.2 = \qty{1305.0}{\celsius}$, \qty{1578}{K}।
\textbf{9.} $T_{\mathrm s}' = 132.9 - 15.3 = \qty{117.6}{K/GPa}$, $T_{\mathrm l}' =
80.0 - 9.6 = \qty{70.4}{K/GPa}$: $(\partial T/\partial P)_F = 117.6 +
0.10\times(70.4 - 117.6) = \qty{112.9}{K/GPa}$; $(\partial T/\partial F)_P =
\qty{314.2}{K}$।
\textbf{10.} $c_p/T = 1000/1578 = 0.634$; अंश $0.634\times112.9 - 12.1 = 71.6
- 12.1 = 59.5$; हर $300 + 0.634\times314.2 = 499.2$: $0.119$, या
\qty{11.9}{\percent} प्रति गीगापास्कल।
\textbf{11.} फ़िट किए गए फलन $F = T'^{1.5}$ के साथ, जहाँ $T'$ सॉलिडस से लिक्विडस तक के मार्ग का अंश है, गलित का पहला प्रतिशत बनाने के लिए
तापमान को $F$ के साथ तीव्रता से बढ़ना पड़ता है, इसलिए सॉलिडस के निकट $(\partial T/\partial F)_P$ बड़ा है
और गलन धीरे शुरू होता है। $F_{\max} = 0.09\times2.66 = 0.24$, \qty{24}{\percent}।
\textbf{12.} $0.13/0.50 = 0.26$: क्लाइनोपाइरॉक्सीन \qty{26}{\percent} गलन पर
समाप्त होता है। \qty{24}{\percent} पर अक्ष के नीचे का अवशेष अब भी थोड़ा धारण करता है:
क्लाइनोपाइरॉक्सीन में निर्धन एक \omterm{def:b2:mantle-melting:peridotites}{लर्ज़ोलाइट}, \omterm{def:b2:mantle-melting:peridotites}{हार्ज़बर्गाइट} की कगार पर।
\textbf{13.} $h_c = 0.24\times82/2 = \qty{9.8}{km}$।
\textbf{14.} $F_{\max}/2 = \qty{12}{\percent}$।
\textbf{15.} \qty{9.8}{km}, \qty{7.1}{km} से \qty{38}{\percent} अधिक है,
सामान्य भूपर्पटी के परास से ऊपर।
\textbf{16.} स्थिर उत्पादकता सबसे गहरे प्रावार को बहुत अधिक गलित देती है;
फ़िट किया गया मॉडल पहला गलित धीरे बनाता है, आरंभ के निकट $F(z)$ कम रहता है, और
$\int F\,\mathrm dz$ छोटा होता है।
\textbf{17.} $F_{\max} = \Pi P_0$ और $H = 30.9P_0$ (km): $h_c = \Pi P_0\times30.9P_0/2
= kP_0^2$, जहाँ $k = 0.09\times30.9/2 = \qty{1.39}{km/GPa^2}$।
\textbf{18.} $P_0 = \sqrt{7.1/1.39} = \sqrt{5.11} = \qty{2.26}{GPa}$।
\textbf{19.} $2.26\times30.9 = \qty{70}{km}$।
\textbf{20.} $T_{\mathrm s}(2.26) = 1085.7 + 300.4 - 26.0 = \qty{1360.0}{\celsius}$,
और $T_p = 1360.0 - 20\times2.26 = 1360.0 - 45.2 = \qty{1315}{\celsius}$।
\textbf{21.} \qty{6.3}{km} के लिए: $P_0 = \sqrt{4.53} = \qty{2.13}{GPa}$, $T_{\mathrm s}
= \qty{1345.6}{\celsius}$, $T_p = \qty{1303}{\celsius}$। \qty{7.9}{km} के लिए: $P_0 =
\sqrt{5.68} = \qty{2.38}{GPa}$, $T_{\mathrm s} = \qty{1373.1}{\celsius}$, $T_p =
\qty{1326}{\celsius}$। भूपर्पटी का प्रकीर्णन $1315 \pm
\qty{12}{\celsius}$ अनुमत करता है।
\textbf{22.} $P_0 = \sqrt{20/1.39} = \sqrt{14.39} = \qty{3.79}{GPa}$; $T_{\mathrm s} =
1085.7 + 503.7 - 73.3 = \qty{1516.1}{\celsius}$; $T_p = 1516.1 - 75.8 =
\qty{1440}{\celsius}$, अर्थात $1440 - 1315 = \qty{125}{\celsius}$ का आधिक्य।
\textbf{23.} दोनों भिन्न साक्ष्यों का प्रयोग करते हैं। लावा सबसे गरम प्रावार को दर्ज करते हैं,
प्रावार-स्तंभ के केंद्र पर; भूपर्पटी उस सारे प्रावार के गलितों का औसत लेती है जो
गलन-क्षेत्र से होकर गुज़रता है, जिसका अधिकांश भाग ठंडा है। सरल मॉडल
उत्पादकता को स्थिर भी रखता है, यद्यपि क्लाइनोपाइरॉक्सीन \qty{26}{\percent} पर समाप्त हो जाता है
जबकि यहाँ $\Pi P_0 = 0.09\times3.79 = 0.34$: वास्तविक गरम प्रावार कम पिघलता है, और
वही भूपर्पटी बनाने के लिए उसे अधिक तापमान चाहिए।
\textbf{24.} $T_{\mathrm s}'(2.26) = 132.9 - 23.1 = \qty{109.8}{K/GPa}$; $\Delta P_0 =
10/(109.8 - 20) = \qty{0.111}{GPa}$; चूँकि $h_c \propto P_0^2$, $\Delta h_c/h_c =
2\times0.111/2.26 = \qty{9.8}{\percent}$, यानी प्रति \qty{10}{\celsius} लगभग दसवाँ भाग।
\textbf{25.} बहुत धीमे कटक पर प्रावार इतनी धीरे उठता है कि वह ठंडे पृष्ठ की ओर
चालन से ऊष्मा खोता है; वह रुद्धोष्म छोड़ देता है, उसका गलन अधिक गहराई पर
रुक जाता है, और त्रिभुज का ऊपरी भाग, जहाँ $F$ सबसे बड़ा है, कट जाता है: भूपर्पटी
पतली होती है।
\textbf{26.} सामान्य कटकों की भूपर्पटी \omterm{def:b2:mantle-melting:tp}{प्रावार का विभव तापमान}
\textbf{\qty[text-series-to-math=true]{1315}{\celsius}} देती है, लगभग
\qty{12}{\celsius} के भीतर: खनिज भौतिकी के \qty{1350}{\celsius} से \qty{35}{\celsius} नीचे,
जो लगभग उतना ही अनिश्चित है, और लावा से मिले 1280 से
\qty{1400}{\celsius} के परास के भीतर। तीन स्वतंत्र तापमापी लगभग
\qty{1300}{\celsius} से \qty{1350}{\celsius} के प्रावार पर सहमत हैं।
\end{solution}
